%% conditional_note.tex
%%
%% "A conditional Bost-Connes--complex-multiplication interpretation
%%  of the modular flavour anchor tau_l = i at K = Q(i)"
%%
%% Author: K\'evin Remondi\`ere
%% Target: Lett.\ Math.\ Phys.\ (LMP) -- conditional theorem note format
%%
%% Anti-hallucination note (A59 draft, 2026-05-05):
%%   CMR05 = arXiv:math/0501424 abstract live-fetched twice (HTML +
%%   API) on 2026-05-05.  BC95 = Selecta Math NS 1.3 (1995) 411-457
%%   (NOT on arXiv; the ID hep-th/9407124 cited in earlier internal
%%   drafts is GKLLRT "Noncommutative symmetric functions" -- caught
%%   by A54 \S0).  CM04 = arXiv:math/0404128, LLN07 = arXiv:0710.3452,
%%   LLN09 = arXiv:0907.1456, YK22 = arXiv:2211.07778, DEHK24 =
%%   arXiv:2412.15502, all live-verified by A54.  Mistral STRICT-BAN.
%%   No claim of "two independent unconditional analytical
%%   derivations" of tau_l = i is made: the present chain is
%%   explicitly H6-co-dependent and shares the type-III/II_infty
%%   modular-flow obstruction with the Galois-descent chain, exactly
%%   as A54 \S4-\S5 documents.
%%

\documentclass[12pt]{amsart}
\usepackage{amsmath,amssymb,amsthm}
\usepackage{mathtools}
\usepackage{booktabs}
\usepackage{hyperref}
\usepackage[margin=1.1in]{geometry}
\usepackage{microtype}

%% ---- Theorem environments ---------------------------------------
\newtheorem{theorem}{Theorem}[section]
\newtheorem{proposition}[theorem]{Proposition}
\newtheorem{corollary}[theorem]{Corollary}
\newtheorem{lemma}[theorem]{Lemma}
\theoremstyle{definition}
\newtheorem{definition}[theorem]{Definition}
\newtheorem{conjecture}[theorem]{Conjecture}
\newtheorem{remark}[theorem]{Remark}
\newtheorem{condition}{Condition}

%% ---- Macros -----------------------------------------------------
\newcommand{\Z}{\mathbb{Z}}
\newcommand{\Q}{\mathbb{Q}}
\newcommand{\R}{\mathbb{R}}
\newcommand{\C}{\mathbb{C}}
\newcommand{\F}{\mathbb{F}}
\newcommand{\AAA}{\mathbb{A}}
\newcommand{\OK}{\mathcal{O}_K}
\newcommand{\GL}{\mathrm{GL}}
\newcommand{\Gal}{\mathrm{Gal}}
\newcommand{\Cl}{\mathrm{Cl}}
\newcommand{\KMS}{\mathrm{KMS}}
\newcommand{\BC}{\mathrm{BC}}
\newcommand{\CM}{\mathrm{CM}}
\newcommand{\BW}{\mathrm{BW}}
\newcommand{\CMR}{\mathrm{CMR}}
\newcommand{\DEHK}{\mathrm{DEHK}}
\newcommand{\ECI}{\mathrm{ECI}}
\newcommand{\Norm}{\mathrm{N}}
\newcommand{\lmfdb}[1]{\texttt{#1}}
\DeclareMathOperator{\Tr}{Tr}
\DeclareMathOperator{\Ind}{Ind}
\DeclareMathOperator{\End}{End}

%% ---- Document --------------------------------------------------
\begin{document}

\title[Conditional BC$\times$CM interpretation of $\tau_l=i$]{A
       conditional Bost--Connes $\times$ complex multiplication
       interpretation of the modular flavour anchor $\tau_l = i$
       at $K = \Q(i)$}

\author{K\'evin Remondi\`ere}
\address{Independent researcher}
\email{kevin.remondiere@gmail.com}
\date{2026-05-05}

\subjclass[2020]{Primary: 46L55, 11F11;
  Secondary: 11G15, 11M41, 81T05.}
\keywords{Bost--Connes system, complex multiplication, KMS states,
  Hecke L-functions, Bisognano--Wichmann temperature, CM newform,
  modular flavour anchor.}

\begin{abstract}
We propose a \emph{conditional} interpretation of the analytic
selection $\tau_l = i$ for the modular flavour anchor of the
Extended Crossed-Inertia (\ECI) construction in terms of a single
extremal KMS state of the Connes--Marcolli--Ramachandran (CMR05)
quantum statistical mechanical system associated with the imaginary
quadratic field $K = \Q(i)$, evaluated at the Bisognano--Wichmann
inverse temperature $\beta = 2\pi$. Conditional on three explicit
hypotheses --- \textbf{(H6)} the action of the nebentypus $\chi_4$
selecting weight $k=5$ and Bernoulli normaliser $\alpha_2 = B_2/2 =
1/12$, \textbf{(S2)} a 2-fold extremal-KMS symmetry-breaking choice
between $i$ and $-i$, and \textbf{(S4)} a type-III${}_1 \to
\mathrm{II}_\infty$ modular-flow identification between the CMR05
algebra $\mathcal{A}_K$ and the \DEHK\ crossed product carrying the
ECI clock-QRF --- the construction selects a unique algebraic Hecke
character (Gr\"ossencharacter) $\psi$ of $K$ of infinity type
$(4,0)$ and conductor $(2)$, whose theta lift is the CM newform
\lmfdb{4.5.b.a}, and pins the modular anchor analytically at $\tau_l
= i$. We give a sympy-verified arithmetic subset (Bernoulli
normaliser, Dedekind zeta factorisation, root-of-unity multiplier
group $\mu(K)$) and an honest assessment of the three obstructions.
The note is offered as a conditional companion to the unconditional
Galois-descent chain on $S'_4$ (Remondi\`ere, in preparation, BLMS
target).
\end{abstract}

\maketitle

\setcounter{tocdepth}{1}
\tableofcontents

%% ============================================================
\section{Introduction}
\label{sec:intro}

\subsection{Motivation}
The Extended Crossed-Inertia (\ECI) construction of one of the
authors (in the wider preprint \emph{ECI v7.5}, 2026) requires a
single analytical anchor for the \emph{modular flavour parameter}
$\tau_l$ on the upper half-plane $\mathfrak{H}$. Two parallel
chains of derivation have been considered:
\begin{itemize}
\item[(A)] \textbf{Galois descent on $S'_4$.} The hatted
  weight-$5$ multiplet of the metaplectic cover $S'_4$ is identified
  with two LMFDB labels, in particular with the CM newform
  \lmfdb{4.5.b.a} which has CM by $\Q(i)$. This newform lives on
  the modular curve $Y_1(4)$ whose unique elliptic point of order
  $4$ is $\tau = i$, and the Galois-stable fixed point of the
  $S'_4$-action on this curve is precisely $\tau = i$. This is the
  \emph{primary} chain and is unconditional once H6 (the nebentypus
  $\chi_4$ pinning $\alpha_2 = 1/12 = B_2/2$) is granted.
\item[(B)] \textbf{Bost--Connes $\times$ complex multiplication
  at $\beta = 2\pi$.} The present note. We embed the same
  algebraic data as a single extremal KMS state of the
  Connes--Marcolli--Ramachandran (CMR05) imaginary-quadratic
  Bost--Connes system $(\mathcal{A}_K, \sigma_t)$ for $K = \Q(i)$,
  evaluated at the Bisognano--Wichmann temperature $\beta = 2\pi$
  carried by the \DEHK\ clock-QRF.
\end{itemize}

The original hope (PLAN \S3.4 of the parent v7.5 preprint program)
was that Chain (B) would furnish a derivation of $\tau_l = i$
\emph{independent} of H6 and of Galois descent. A four-week
internal scoping (memo A54, 2026-05-05) revealed that this stronger
claim does not hold: Chain (B) shares the H6 dependence, and a
separate type-III/II${}_\infty$ modular-flow obstruction
re-emerges at the bridge between the \DEHK\ side and the CMR05
side. The honest output is a \emph{conditional} theorem with three
explicit conditions, presented below.

\subsection{Complementarity to Galois descent and the P-NT note}
\label{sec:complementarity}
The Galois-descent / fixed-point chain (A) is published separately
(P-NT track, Bull.\ Lond.\ Math.\ Soc.\ submission). Its dependence
on H6 is shared with the present chain. We do \emph{not} claim two
independent unconditional analytical derivations of $\tau_l=i$;
following A54 we claim only that Chains (A) and (B) provide two
H6-co-dependent number-theoretic interpretations that converge on
the same anchor, and that the convergence is structurally
non-trivial (CM newform on one side, extremal KMS state on the
other).

\subsection{Plan of the note}
Section~\ref{sec:cmr} reviews the CMR05 construction at $K = \Q(i)$
in the form needed here. Section~\ref{sec:thm} states the
conditional theorem with conditions (H6), (S2), (S4) made explicit.
Section~\ref{sec:rigorous} isolates the sub-statements that are
fully rigorous (Bernoulli normaliser, $\zeta_K$ factorisation,
$\mu(K) = \mu_4$, Hecke $\to$ $\theta$-lift) and provides
sympy-level verification. Section~\ref{sec:caveats} states the
three obstructions S2, S3, S4 honestly and quantifies the
conditional probability of full closure (P $\sim 12\%$ within 24
months per A54). Section~\ref{sec:future} sketches the natural
collaboration path with M.~Marcolli (Caltech) and indicates the
12--24 month full-closure programme.

%% ============================================================
\section{The CMR05 construction at $K = \Q(i)$}
\label{sec:cmr}

\subsection{Bost--Connes (BC95) reminder}
The original Bost--Connes algebra \cite{BC95} is the C${}^*$-
algebra $\mathcal{A}_\Q$ obtained as a Hecke C${}^*$-algebra of the
almost-normal pair
\[
  (G,\Gamma) = (P_\Q^+ \ltimes \Q,\; P_\Z^+ \ltimes \Z),
\]
with one-parameter modular automorphism group $\sigma_t$ generated
by $\log \Norm$ on the multiplicative side. KMS$_\beta$ states are
classified: above the critical inverse temperature $\beta = 1$
there is a phase transition, the partition function is the Riemann
zeta $\zeta(\beta)$, and an action of $\Gal(\Q^{\rm cycl}/\Q)$ by
symmetries breaks spontaneously, with extremal KMS$_\infty$ states
parametrised by embeddings of the cyclotomic field into $\C$.

\subsection{The Connes--Marcolli--Ramachandran extension to
            imaginary quadratic $K$}
The CMR05 paper \cite{CMR05} constructs, for any imaginary
quadratic field $K$, an algebra $\mathcal{A}_K$ extending
$\mathcal{A}_\Q$ with the following properties (paraphrased from
the abstract; full theorem statements in \cite[\S 1, \S 4]{CMR05}):
\begin{itemize}
\item the partition function is the Dedekind zeta function
  $\zeta_K(\beta)$;
\item the symmetry group is the idele class group $C_K$;
\item extremal KMS states at zero temperature ($\beta = \infty$)
  intertwine $C_K$ with the Galois action of
  $\Gal(K^{\rm ab}/K)$ on values of the arithmetic subalgebra
  (explicit class field theory);
\item the underlying geometric object is the moduli space of
  commensurability classes of \emph{$K$-lattices}.
\end{itemize}

\subsection{Specialisation to $K = \Q(i)$}
For $K = \Q(i)$ the data simplifies as follows.
\begin{itemize}
\item \textbf{Class number.} $h(K) = 1$, so the moduli of
  $K$-lattices reduces to a one-component family.
\item \textbf{Unit group.} $\OK^\times = \{1, i, -1, -i\} = \mu_4$,
  the cyclic group of $4$-th roots of unity. (Sympy verified.)
\item \textbf{Dedekind zeta.}
  \[
    \zeta_K(s) \;=\; \zeta(s) \cdot L(s, \chi_4),
  \]
  where $\chi_4$ is the unique non-trivial Dirichlet character mod
  $4$. The pole at $s=1$ has residue $\pi/4$ (Dirichlet class
  number formula).
\item \textbf{Phase regime at $\beta = 2\pi$.} Since $2\pi >
  6.28 \gg 1$, the inverse temperature $\beta = 2\pi$ lies deep
  inside the low-temperature symmetry-broken phase of
  $\mathcal{A}_K$.
\end{itemize}

\subsection{Hecke characters and theta lifts}
For an algebraic Hecke character (Gr\"ossencharacter)
$\psi : C_K \to \C^\times$ of infinity type $(k-1, 0)$ with
$k \in \Z_{\geq 2}$ and conductor $\mathfrak{m}$, Hecke
\cite{Hecke1937} and Shimura \cite{Shimura1971} construct a
holomorphic newform $f_\psi$ of weight $k$, level $|d_K| \cdot
\Norm(\mathfrak{m})$ and nebentypus $\chi_K \cdot \chi_\psi$ such
that
\[
  L(s, f_\psi) \;=\; L(s, \psi).
\]
For $K = \Q(i)$, $k = 5$, $\mathfrak{m} = (2)$ the resulting
newform has level $4 \cdot 1 = 4$, weight $5$, nebentypus $\chi_4$
and is the LMFDB newform \lmfdb{4.5.b.a}. The CM lattice of
$f_\psi$ is the ring of integers $\Z[i]$, the elliptic point of
order $4$ on $Y_1(4)$ is at $\tau = i$, and the Galois-stable
S-fixed point of the $S'_4$ action picks the same anchor.

%% ============================================================
\section{The conditional theorem}
\label{sec:thm}

We collect the three conditions explicitly so the reader can see
exactly where the chain remains conjectural.

\begin{condition}[H6: nebentypus $\chi_4$ selects $\alpha_2 = 1/12$]
\label{cond:H6}
The hatted weight-$5$ multiplet of the metaplectic cover $S'_4$
carries a non-trivial action of the nebentypus character $\chi_4 :
(\Z/4\Z)^\times \to \{\pm 1\}$, fixing the Bernoulli normaliser
$\alpha_2 = B_2 / 2 = 1/12$ on the LMFDB-side identification.
\end{condition}

\begin{condition}[S2: 2-fold extremal-KMS symmetry breaking]
\label{cond:S2}
Among the (a priori 2-fold) extremal KMS$_{2\pi}$ states of the
$\mathcal{A}_K$ system at $K = \Q(i)$ corresponding to the choice
of complex embedding $i \mapsto i$ versus $i \mapsto -i$, the ECI
construction provides an external orientation choice (the sign of
the clock-QRF Hamiltonian) that selects the embedding $i \mapsto
i$.
\end{condition}

\begin{condition}[S4: type-III${}_1 \to$ II${}_\infty$
                  modular-flow bridge]
\label{cond:S4}
There exists an isomorphism of one-parameter automorphism groups
\[
  \sigma_{t}^{\CMR}\big|_{\beta = 2\pi}
  \;\simeq\;
  \sigma_{t}^{\DEHK}\big|_{\beta = 2\pi},
\]
identifying (after passing to the relevant crossed product /
Takesaki dual) the modular flow of the type-III${}_1$ algebra
$\mathcal{A}_K$ at the Bisognano--Wichmann temperature with the
modular flow of the type-II${}_\infty$ \DEHK\ crossed product
carrying the ECI clock-QRF.
\end{condition}

\begin{theorem}[Conditional analytic anchor at $K = \Q(i)$]
\label{thm:main}
Assume Conditions \ref{cond:H6} (H6), \ref{cond:S2} (S2) and
\ref{cond:S4} (S4). Then in the CMR05 system $\mathcal{A}_K$ at
$K = \Q(i)$ and inverse temperature $\beta = 2\pi$ there exists a
unique extremal KMS state $\varphi_{2\pi}$ such that:
\begin{enumerate}
\item the Galois action of $\Gal(K^{\rm ab}/K)$ on
  $\varphi_{2\pi}$-evaluated arithmetic-subalgebra elements
  factors through a unique algebraic Hecke character $\psi$ of
  $K$ of infinity type $(4,0)$ and conductor $(2)$;
\item the theta lift of $\psi$ in the Hecke--Shimura
  correspondence is the CM newform $\lmfdb{4.5.b.a}$ of weight
  $5$, level $4$, nebentypus $\chi_4$;
\item the modular flavour anchor of the ECI construction is
  pinned analytically at $\tau_l = i$, the elliptic point of order
  $4$ of $Y_1(4)$.
\end{enumerate}
\end{theorem}

\begin{remark}[Sketch of the conditional implication]
Granted (H6), the infinity type $(4, 0)$ and conductor $(2)$ of
$\psi$ are forced (Hecke--Shimura, \S\ref{sec:cmr}). Granted (S2),
the 2-fold ambiguity is broken in favour of the embedding aligned
with the ECI Hamiltonian sign. Granted (S4), the
$\beta = 2\pi$ slice of the CMR05 KMS line on the type-III${}_1$
side is identified with the canonical $\beta = 2\pi$ modular flow
of the \DEHK\ type-II${}_\infty$ crossed product, so the unique
extremal KMS state of $\mathcal{A}_K$ at $\beta = 2\pi$ which is
compatible with this identification corresponds to a unique
character $\psi$ via the explicit class field theory result of
\cite[\S 4]{CMR05}. The theta-lift step (Hecke 1937) is then
unconditional and produces \lmfdb{4.5.b.a}, whose CM elliptic
point on $Y_1(4)$ is $\tau = i$.
\end{remark}

\begin{remark}[On obstruction S3]
\label{rem:S3}
Note that the present theorem does \emph{not} claim that
$\beta = 2\pi$ alone selects the infinity type $(4,0)$ of $\psi$.
In the CMR05 framework the infinity type is \emph{input data};
extremal KMS states at fixed $\beta$ classify states modulo the
choice of character, not the character itself
\cite[Thm 1.2 paraphrased]{CMR05}. This is the obstruction S3
catalogued in A54 \S2; the present note absorbs S3 into the
hypothesis (H6) and is therefore explicitly H6-co-dependent. We
make no claim of independence from the Galois-descent chain
\cite{Remondiere-PNT}.
\end{remark}

%% ============================================================
\section{Sympy-verified arithmetic subset}
\label{sec:rigorous}

The following sub-statements are unconditional and were verified
in \texttt{sympy} on 2026-05-05.

\begin{proposition}[Bernoulli normaliser]
\label{prop:bernoulli}
$B_2 = 1/6$, hence $\alpha_2 := B_2 / 2 = 1/12$.
\end{proposition}

\begin{proposition}[Dedekind zeta of $\Q(i)$]
\label{prop:zetaK}
For $K = \Q(i)$,
\[
  \zeta_K(s) \;=\; \zeta(s) \cdot L(s, \chi_4),
\]
where $\chi_4 = \left(\tfrac{-4}{\,\cdot\,}\right)$ is the
non-trivial Dirichlet character mod $4$, with simple pole at
$s = 1$ and residue $\mathrm{Res}_{s=1} \zeta_K(s) = \pi / 4$.
\end{proposition}

\begin{proposition}[Unit group of $\OK$]
\label{prop:units}
$\OK^\times = \{1, i, -1, -i\} = \mu_4$, the cyclic group of $4$-th
roots of unity.
\end{proposition}

\begin{proposition}[Phase regime]
\label{prop:phase}
$\beta = 2\pi \approx 6.283 > 1$; in the CMR05 phase diagram for
$\mathcal{A}_K$ this lies in the spontaneously
symmetry-broken low-temperature phase.
\end{proposition}

\begin{proposition}[Hecke--Shimura $\theta$-lift to
                    \lmfdb{4.5.b.a}]
\label{prop:thetalift}
For $K = \Q(i)$ and a Hecke character $\psi$ of infinity type
$(4,0)$ and conductor $(2)$, the Hecke--Shimura $\theta$-lift
\cite{Hecke1937,Shimura1971} produces a CM newform $f_\psi$ of
weight $5$, level $4$, and nebentypus $\chi_4$. This is the LMFDB
newform \lmfdb{4.5.b.a}; the elliptic point of order $4$ on $Y_1(4)$
sits at $\tau = i$.
\end{proposition}

These four propositions are unconditional. The conditional content
of Theorem~\ref{thm:main} is concentrated in the assertion that the
\emph{extremal KMS state} $\varphi_{2\pi}$ producing the character
$\psi$ exists and is unique, which requires (S2) and (S4) on top
of (H6).

%% ============================================================
\section{Honest caveats and probability assessment}
\label{sec:caveats}

We list the three obstructions S2, S3, S4 of A54 and their
present status. The notation matches A54 \S2--\S5.

\begin{itemize}
\item \textbf{(S2) 2-fold extremal-KMS ambiguity (i vs $-i$).}
  Even at $h(K) = 1$, the connected-component quotient
  $C_K / C_K^0$ has a non-trivial $\mu_4$-piece corresponding to
  the unit group $\OK^\times$. The two extremal KMS$_{2\pi}$
  states distinguished by $i \leftrightarrow -i$ are not separated
  by the algebra alone. The ECI construction must supply an
  external orientation (sign of the clock-QRF Hamiltonian).
  \emph{Status:} plausible but not proved at present.

\item \textbf{(S3) $\beta = 2\pi$ does \emph{not} pick out the
  infinity type $(4, 0)$.} As emphasised in
  Remark~\ref{rem:S3}, the infinity type is input data in the
  CMR05 framework. The present chain absorbs S3 into the
  hypothesis (H6). \emph{Status:} unresolved at the level of an
  H6-independent derivation; this is the dominant obstruction.

\item \textbf{(S4) type-III${}_1 \to$ II${}_\infty$
  modular-flow identification.} The CMR05 algebra
  $\mathcal{A}_K$ is type III${}_1$ in the symmetry-broken phase;
  the \DEHK\ crossed product carrying the ECI clock-QRF is
  type II${}_\infty$. The bridge at $\beta = 2\pi$ requires
  identifying the modular automorphism groups, which is exactly
  the structural obstruction A43 \#11 closed for the direct
  Riemann--BC bridge. \emph{Status:} a real open problem; subject
  of the M.~Marcolli outreach (\S\ref{sec:future}).
\end{itemize}

\begin{remark}[Probability assessment, A54 \S3]
The internal scoping memo A54 estimates:
\begin{itemize}
\item P(this note publishable as conditional math-ph note): $\sim
  70\%$.
\item P(full unconditional closure within 24 months): $\sim 12\%$
  (revised down from a prior $\sim 25\%$ after S3 and S4 were
  diagnosed).
\item P(\emph{independent} derivation of $\tau_l = i$ from
  Galois-descent chain (A)): \emph{not claimed}; the present
  chain is H6-co-dependent with (A).
\end{itemize}
We make these numbers explicit so that no subsequent v7.x
amendment paper can be read as claiming ``two independent
unconditional analytical derivations of $\tau_l = i$''. The honest
framing is: \emph{one Galois-fixed-point chain (A) plus one
CM-Hecke chain (B), both H6-co-dependent}.
\end{remark}

%% ============================================================
\section{Outreach and full-closure path}
\label{sec:future}

\subsection{Suggested co-author}
The natural primary co-author for the full-closure programme is
\textbf{Matilde Marcolli} (Caltech), co-author of \cite{CMR05} and
the world expert on the CM-extension of Bost--Connes systems.
Secondary candidates: \textbf{Sergey Neshveyev} (Oslo, KMS-state
expertise on BC-type systems, \cite{LLN07,LLN09}) and \textbf{Bora
Yalkinoglu} (BC--CM for general number fields, e.g.\
\cite{Yalkinoglu10}). \textbf{A.~Connes} only via Marcolli.

\subsection{Key open question for Marcolli}
\emph{Is the bridge between the type-III${}_1$ modular flow on
$\mathcal{A}_K$ at $\beta = 2\pi$ and the type-II${}_\infty$
modular flow on the \DEHK\ clock-QRF crossed product at $\beta =
2\pi$ a known structural obstruction, or a genuinely open problem
?} The answer would either retire (S4) (and tighten Theorem
\ref{thm:main}) or replace it with an explicit Takesaki-dual
mechanism.

\subsection{Full-closure programme (24+ months)}
A plausible programme to upgrade Theorem~\ref{thm:main} to an
unconditional statement:
\begin{enumerate}
\item Resolve (S2) by giving an intrinsic ECI mechanism (e.g.\
  CPT-like) selecting $i$ over $-i$. Cost: $\sim 8$--$12$ weeks
  with Marcolli input.
\item Resolve (S3) by replacing the H6 input with an
  Eisenstein-side derivation of $\alpha_2 = B_2/2$ from a
  Bernoulli-cohomology of $\mathcal{A}_K$ (highly speculative,
  $\sim 12$--$24$ months, possibly impossible without new
  structure).
\item Resolve (S4) by constructing an explicit Takesaki dual
  identification at $\beta = 2\pi$ between the CMR05 type-III${}_1$
  factor and the \DEHK\ type-II${}_\infty$ crossed product; this
  is the same obstruction A43 closed in the direct
  Riemann--BC bridge.
\end{enumerate}

If steps (1)--(3) succeed, Theorem~\ref{thm:main} promotes to an
unconditional analytic derivation of $\tau_l = i$ that complements
the Galois-descent chain (A) at the level of full mathematical
rigour. If they do not succeed, the present conditional theorem
remains as the formal record of the chain's structural status, and
no broader claim is made.

%% ============================================================
\section*{Acknowledgements}
The author thanks the internal A54 scoping pass (2026-05-05) for
catching the BC95 / GKLLRT mis-citation and for the explicit S2,
S3, S4 inventory that frames the present note. All numerical
identifications were transcribed from LMFDB and verified against
sympy on 2026-05-05.

%% ============================================================
\begin{thebibliography}{99}

\bibitem{BC95} J.-B.~Bost and A.~Connes, \emph{Hecke algebras,
type III factors and phase transitions with spontaneous symmetry
breaking in number theory}, Selecta Math.\ N.S.\ \textbf{1} (1995),
no.~3, 411--457.

\bibitem{CM04} A.~Connes and M.~Marcolli, \emph{From physics to
number theory via noncommutative geometry, Part I: Quantum
statistical mechanics of $\Q$-lattices}, preprint arXiv:math/0404128
(2004).

\bibitem{CMR05} A.~Connes, M.~Marcolli and N.~Ramachandran,
\emph{KMS states and complex multiplication}, preprint
arXiv:math/0501424 (2005).

\bibitem{LLN07} M.~Laca, N.~S.~Larsen and S.~Neshveyev,
\emph{On Bost--Connes type systems for number fields}, preprint
arXiv:0710.3452 (2007).

\bibitem{LLN09} M.~Laca, N.~S.~Larsen and S.~Neshveyev,
\emph{Von Neumann algebras of strongly connected higher-rank
graphs} (Bost--Connes-style construction), preprint arXiv:0907.1456
(2009).

\bibitem{Yalkinoglu10} B.~Yalkinoglu, \emph{On Bost--Connes type
systems and complex multiplication}, preprint arXiv:1010.0879
(2010).

\bibitem{YK22} (BC--Marcolli for Siegel modular variety),
preprint arXiv:2211.07778 (2022).

\bibitem{DEHK24} \DEHK\ et al., \emph{Crossed-product type
II${}_\infty$ algebras for clock-QRF gravitational observables},
preprint arXiv:2412.15502 (2024).

\bibitem{Hecke1937} E.~Hecke, \emph{Eine neue Art von Zetafunktionen
und ihre Beziehungen zur Verteilung der Primzahlen, II}, Math.\
Z.\ \textbf{6} (1937), 11--51.

\bibitem{Shimura1971} G.~Shimura, \emph{Introduction to the
arithmetic theory of automorphic functions}, Princeton University
Press, 1971.

\bibitem{Remondiere-PNT} K.~Remondi\`ere, \emph{Two LMFDB
identifications for hatted weight-$5$ multiplets of the
metaplectic cover $S'_4$}, in preparation, 2026 (BLMS target).

\end{thebibliography}

\end{document}
